\subsection*{Part 6, Mission: Make your own block}
\textbf{Triangle:}
\begin{codenonum}
def triangle(size):
    for i in range(3):
        forward(size)
        right(120)
\end{codenonum}
\textbf{Colour input:}
\begin{codenonum}
def square(size, colour):
    pencolor(colour)
    for i in range(4):
        forward(size)
        right(90)
\end{codenonum}
\textbf{No brackets:} a line with just \py|square| does nothing, and there is no error. Without the brackets you only say the block's name. You must call it, for example \py|square(60)|.

\subsection*{Part 6, Mission: Shape machine}
\py|flower(8, 4, 60)| draws 8 petals, and each petal is a square. With \py|animationOff()| the turtle jumps instead of walking, so the drawing appears faster. Colour input:
\begin{codenonum}
def flower(petals, sides, size, colour):
    pencolor(colour)
    for i in range(petals):
        polygon(sides, size)
        right(360 / petals)
\end{codenonum}
Call it with four inputs, for example \py|flower(12, 6, 40, "red")|.

\subsection*{Part 6, Mission: Answer machine}
\begin{codenonum}
def half(n):
    return n / 2

def perimeter(width, height):
    return 2 * width + 2 * height

def is_even(n):
    return n % 2 == 0
\end{codenonum}
\py|print(half(9))| shows \texttt{4.5}. \py|is_even(7)| is \py|False|, so the \py|if| part is skipped.

With \py|print(n * 2)| in line 2, \py|double| only \textbf{shows} the number and gives back \texttt{None}. So line 7 shows \texttt{8} and then \texttt{None}. Line 11 then tries \py|area(None, 10)|, and Python stops with a \texttt{TypeError}, because it can't multiply nothing by 10.

\subsection*{Part 6, Mission: Sprite army}
\begin{codenonum}
def turn_all(angle):
    for s in army:
        s.left(angle)
\end{codenonum}
Call \py|turn_all(90)| before marching, and the army marches up. For the race, put \py|import random| at the top and use \py|s.forward(random.randint(10, 80))| in \texttt{march}. \py|army[0]| is the first sprite made: the red one at the top (lists count from 0).

\subsection*{Part 6, Translate it yourself}
\begin{codenonum}
from ColabTurtlePlus.Turtle import *
clearscreen()
initializeTurtle((480, 360))

def star(size):
    for i in range(5):
        forward(size)
        right(144)

pencolor("gold")
star(100)
star(50)
\end{codenonum}
You should see a big gold star with a small one inside it, sharing one point.

\subsection*{Part 6, Two-line puzzle}
\begin{codenonum}
def triple(n):
    return n * 3
\end{codenonum}
\py|triple(2)| is 6, and \py|triple(6)| is 18.
